\section*{अध्याय \ref{ch:b2:population-genetics} --- समष्टि आनुवंशिकी}

\begin{solution}{exo:b2:population-genetics:1}
$p(M) = (1280 + 320)/2000 = 0.8$, $q(N) = 0.2$। अपेक्षित $MM$:
$0.64\times 1000 = 640$; $MN$: $2\times 0.8\times 0.2\times 1000 =
320$; $NN$: 40 --- यानी ठीक वही संख्याएँ जो देखी गईं। वह समष्टि इस स्थल पर
साम्य में है।
\end{solution}

\begin{solution}{exo:b2:population-genetics:2}
$q = \sqrt{1/20000} = 0.0071$; वाहक $2pq \approx 0.014$, यानी 70 में
से एक व्यक्ति। वाहक प्रभावित लोगों से $2p/q \approx
280$ गुना अधिक हैं।
\end{solution}

\begin{solution}{exo:b2:population-genetics:3}
\omterm{def:b2:population-genetics:fitness}{अनुकूलता}: किसी \omterm{def:b2:meiosis-heredity:vocabulary}{जीन-प्ररूप} की संतति की प्रत्याशित संख्या,
सर्वोत्तम के सापेक्ष। \omterm{def:b2:population-genetics:fitness}{वरण गुणांक} $s$: यानी \omterm{def:b2:population-genetics:fitness}{अनुकूलता} की कमी, $w = 1 -
s$। वंशागतिता $h^{2}$: उस लक्षणप्ररूपी प्रसरण का वह अंश
जो योगात्मक आनुवंशिक है, यानी समतुल्य रूप से जनकों पर संतति का
ढाल। बल: \omterm{def:b2:genome-diversification:mutation}{उत्परिवर्तन}, वरण, अपवाह, प्रवास, अयादृच्छिक
प्रजनन (अंतिम \omterm{def:b2:meiosis-heredity:vocabulary}{जीन-प्ररूप} बारंबारताएँ बदलता है, \omterm{def:b2:population-genetics:frequencies}{युग्मविकल्पी बारंबारताएँ} नहीं)।
\end{solution}

\begin{solution}{exo:b2:population-genetics:4}
दिशात्मक: मिर्चदार पतंगे में कालापन, जहाँ कारक वे चिड़ियाँ हैं जो
कालिख से काले तनों पर देखकर शिकार करती हैं। स्थायीकारी: मानव जन्म-भार,
जहाँ कारक बहुत छोटे तथा बहुत बड़े शिशुओं की मृत्यु है। संतुलनकारी:
हँसिया-कोशिका, जहाँ कारक मलेरिया ($AA$ के विरुद्ध) तथा रक्ताल्पता ($SS$ के
विरुद्ध) साथ-साथ हैं।
\end{solution}

\begin{solution}{exo:b2:population-genetics:5}
$q_t = q_0/(1 + tq_0)$: आधा करने में $1/q_0 = 50$ पीढ़ियाँ लगती हैं;
और $0.001$ तक पहुँचने में $1/0.001 - 1/0.02 = 950$ पीढ़ियाँ। उस
\omterm{def:b2:meiosis-heredity:vocabulary}{युग्मविकल्पी} की लगभग सारी प्रतियाँ विषमयुग्मजों में बैठी हैं, यानी वरण को
अदृश्य; और प्रभावितों को प्रजनन से रोकना (जो प्रकृति पहले ही कर देती है,
क्योंकि वे मर चुके हैं) किसी मानव जीवनकाल में मापने लायक कुछ भी नहीं
बदलता।
\end{solution}

\begin{solution}{exo:b2:population-genetics:6}
$p = 0.3$: $\bar w = 0.09 + 0.42 + 0.8\times 0.49 = 0.902$; $p' =
(0.09 + 0.21)/0.902 = 0.333$; $\Delta p = +0.033$। $p = 0.9$: $\bar w
= 0.998$, $p' = 0.9018$, $\Delta p = +0.0018$। $p = 0.9$ पर उस समष्टि का केवल
$q^{2} = 0.01$ $aa$ है: जिस \omterm{def:b2:meiosis-heredity:vocabulary}{युग्मविकल्पी} के विरुद्ध वरण काम करता है
वह विषमयुग्मजों में छिपा है, और $\Delta p = spq^{2}/\bar w$
$q^{2}$ का गुणक ढोता है।
\end{solution}

\begin{solution}{exo:b2:population-genetics:7}
$s = 0.12$, जो $AA$ के विरुद्ध है, और $t = 0.8$, जो $SS$ के विरुद्ध, के साथ $\hat q =
s/(s + t) = 0.12/0.92 = 0.13$। प्रभावित जन्म $\hat q^{2} = 0.017$,
यानी लगभग 60 में एक बच्चा --- यानी उस रक्षा का दाम जो वे
विषमयुग्मज भोगते हैं।
\end{solution}

\begin{solution}{exo:b2:population-genetics:8}
\omterm{def:b2:meiosis-heredity:vocabulary}{अप्रभावी} घातक: $\hat q = \sqrt{\mu} = \sqrt{2\times 10^{-5}} =
0.0045$; प्रभावित जन्म $\hat q^{2} = \mu = 2\times 10^{-5}$, यानी
\num{50000} में एक। \omterm{def:b2:meiosis-heredity:vocabulary}{प्रभावी} घातक: हर प्रति उसी पीढ़ी में हटा दी जाती है
जिसमें वह प्रकट हो, इसलिए प्रभावित जन्म केवल वे नए
\omterm{def:b2:genome-diversification:mutation}{उत्परिवर्तन} हैं, $2\mu = 4\times 10^{-5}$ (प्रति जन्म दो युग्मक), यानी
\num{25000} में एक --- वह \omterm{def:b2:meiosis-heredity:vocabulary}{अप्रभावी} अपने \omterm{def:b2:meiosis-heredity:vocabulary}{युग्मविकल्पी} दो सौ गुना छिपाता है,
और वह \omterm{def:b2:meiosis-heredity:vocabulary}{प्रभावी} कुछ भी नहीं छिपाता।
\end{solution}

\begin{solution}{exo:b2:population-genetics:9}
$H_t = 0.5\,(1 - 1/100)^{t}$: 10 पीढ़ियों बाद $0.45$; 50 बाद
$0.30$; और 200 बाद $0.067$। $(1 - 1/2N)^{100} = 0.9$ के लिए:
$100/(2N) \approx -\ln 0.9 = 0.105$, इसलिए $N \approx 475$, यानी लगभग 500
प्रजनन करते व्यष्टि।
\end{solution}

\begin{solution}{exo:b2:population-genetics:10}
बीस जीन-प्रतियाँ: $P(\text{अनुपस्थित}) = 0.9^{20} = 0.12$। 20 में से कोई एक
प्रति $0.05$ बारंबारता पर है, और कोई उदासीन \omterm{def:b2:meiosis-heredity:vocabulary}{युग्मविकल्पी} अपनी बारंबारता के बराबर
प्रायिकता से नियत होता है: यानी $0.05$। मान लीजिए
\num{100000} पक्षियों की किसी मुख्यभूमि पर किसी एक प्रति की नियत होने की प्रायिकता $5\times 10^{-6}$
है। वह द्वीप ऐसी जगह है जहाँ दुर्लभ \omterm{def:b2:meiosis-heredity:vocabulary}{युग्मविकल्पी} खोते हैं और
जहाँ जो बचें वे सब कुछ ले सकते हैं: अपवाह दरिद्र भी करता है और
विभेदित भी।
\end{solution}

\begin{solution}{exo:b2:population-genetics:11}
$R = h^{2}S = 0.3\times 1.4\times \qty{1000}{L} = \qty{420}{L}$ प्रति
पीढ़ी। ऊपर के \qty{1}{\%} से साँड़ों के साथ, औसत
अंतर $(1.4 + 2.7)/2 = 2.05$ मानक विचलन है और $R =
0.3\times 2050 = \qty{615}{L}$। वह अनुक्रिया इसलिए धीमी पड़ती है कि वरण
उस योगात्मक प्रसरण को चुका देता है --- अनुकूल \omterm{def:b2:meiosis-heredity:vocabulary}{युग्मविकल्पी} नियतन तक चले जाते हैं
और $h^{2}$ गिरती है --- और इसलिए भी कि दूसरे लक्षणों में सहसंबद्ध परिवर्तन
(जननक्षमता, स्वास्थ्य) ऐसे दाम लगाते हैं जिन्हें केवल उपज पर वरण
देखता ही नहीं।
\end{solution}

\begin{solution}{exo:b2:population-genetics:12}
वरण \omterm{def:b2:meiosis-heredity:vocabulary}{लक्षण-प्ररूप} देखता है। किसी विषमयुग्मज में कोई \omterm{def:b2:meiosis-heredity:vocabulary}{अप्रभावी} \omterm{def:b2:meiosis-heredity:vocabulary}{युग्मविकल्पी}
कोई \omterm{def:b2:meiosis-heredity:vocabulary}{लक्षण-प्ररूप} नहीं बनाता, इसलिए वरण उसे हटा नहीं सकता, और उसकी
बारंबारता $1/t$ की तरह गिरती है, धीरे और धीरे। कोई उदासीन \omterm{def:b2:genome-diversification:mutation}{उत्परिवर्तन}
कोई अनुकूलता-अंतर नहीं बनाता; उसका भाग्य केवल अपवाह का है, और वह
$1/2N$ प्रायिकता से नियत होता है। शून्य वंशागतिता वाला कोई लक्षण
ज़ोर से वरा जा सकता है --- बचे हुए उस समष्टि से बहुत भिन्न हो
सकते हैं --- फिर भी अगली पीढ़ी अपरिवर्तित रहती है, क्योंकि
वे भेद वंशागत नहीं थे। वरण केवल ऐसे वंशागत भेदों पर काम कर
सकता है जो \omterm{def:b2:population-genetics:fitness}{अनुकूलता} में अंतर लाएँ; और बाक़ी सब
संयोग से विकसित होता है या बिलकुल नहीं।
\end{solution}

\begin{solution}{pb:b2:population-genetics:1}
\textbf{1.} काला अंश $1 - q^{2} = 1 - 0.999^{2} = 0.002$,
यानी 500 में एक पतंगा; और विषमयुग्मजों में $M$ युग्मविकल्पियों का अंश
$2pq/(2pq + 2p^{2}) = q = 0.999$ है: यानी मूलतः सारे।
\textbf{2.} $\bar w = 1 - sq^{2}$ के साथ, वरण के बाद पीले \omterm{def:b2:meiosis-heredity:vocabulary}{युग्मविकल्पी} की
प्रतियाँ $pq\cdot 1 + q^{2}(1 - s)$ हैं, $\bar w$ में से, इसलिए $q'
= (pq + q^{2} - sq^{2})/(1 - sq^{2}) = q(1 - sq)/(1 - sq^{2})$।
\textbf{3.} $q_1 = 0.999\times(1 - 0.2997)/(1 - 0.2994) = 0.99857$:
यानी $-0.0014$ का परिवर्तन।
\textbf{4.} $q = \sqrt{0.1} = 0.32$।
\textbf{5.} वह पुनरावर्तन, दोहराया जाए तो, $s =
0.3$ के लिए 28 पीढ़ियाँ लेता है; और वह सन्निकटन वही कोटि देता है (वह परिवर्तन
पहले $-0.0003$ है, फिर $p$ के बढ़ने के साथ तेज़ होता है)। देखी गई 50
पीढ़ियाँ किसी छोटी $s$ से मेल खाती हैं, लगभग $0.2$ (वह पुनरावर्तन
$s = 0.2$ के लिए 45 पीढ़ियाँ देता है): यानी प्रति पीढ़ी \qtyrange{20}{30}{\%} का वरण,
जो अधिकांश युग्मविकल्पियों के मानकों से विशाल है।
\textbf{6.} $w_{MM} = w_{Mm} = 1 - s$, $w_{mm} = 1$ के साथ: $\bar w = 1
- s(1 - q^{2})$ और $p' = p(1 - s)/\bar w$, इसलिए $\Delta p = p[(1 - s) -
\bar w]/\bar w = -spq^{2}/\bar w$। वह पीला रूप केवल
समयुग्मज $mm$ के रूप में अनुकूल है, यानी आरंभ में कोई अंश $q^{2} = 0.1$: इसलिए वह
वापसी धीरे शुरू होती है। ज्यों-ज्यों $q \to 1$, त्यों-त्यों $\Delta p \approx -sp$ और $M$
चरघातांकी रूप से घटता है, हर पीढ़ी $1 - s$ के गुणक से: किसी
\omterm{def:b2:meiosis-heredity:vocabulary}{प्रभावी} \omterm{def:b2:meiosis-heredity:vocabulary}{युग्मविकल्पी} के पास छिपने को कहीं नहीं है। वह चढ़ाई उसका दर्पण-प्रतिबिंब थी:
$\Delta p \approx +sp$ जब तक $M$ दुर्लभ था (यानी चरघातांकी आरंभ), फिर
$\Delta q \approx -sq^{2}p \to$ कोई रेंगन, ज्यों ही वह पीला \omterm{def:b2:meiosis-heredity:vocabulary}{युग्मविकल्पी}
विषमयुग्मजों में छिपा। वह पुनरावर्तन
$p = 0.68$ से $0.001$ तक वापसी के लिए 27 पीढ़ियाँ देता है, जिनमें से 16 उस \omterm{def:b2:meiosis-heredity:vocabulary}{युग्मविकल्पी} को \qty{5}{\%} तक लाने में।
\textbf{7.} $24$ जीन-प्रतियाँ: $0.95^{24} = 0.29$।
\textbf{8.} $H_1 = 0.4\,(1 - 1/24) = 0.383$।
\textbf{9.} $H_{100} = 0.383\,(1 - 1/80)^{100} = 0.383\times 0.284 =
0.109$: यानी मुख्यभूमि के $0.4$ का \qty{27}{\%}।
\textbf{10.} $P(\text{नियतन}) = 1/(2N) = 1/80 = 0.0125$; नियतन तक
प्रत्याशित काल $4N = 160$ पीढ़ियाँ।
\textbf{11.} अंतःप्रजनन के बिना, $q^{2} = 0.0025$; और $F = 0.3$ के साथ,
$q^{2} + Fpq = 0.0025 + 0.3\times 0.95\times 0.05 = 0.017$: यानी सात
गुना अधिक प्रभावित पक्षी।
\textbf{12.} प्रति पीढ़ी एक प्रवासी ($m = 1/40$) उस द्वीप को उदासीन
स्थलों पर विचलित होने से रोकने और अपवाह ने जो \omterm{def:b2:meiosis-heredity:vocabulary}{युग्मविकल्पी}
खोए उन्हें फिर भरने भर को काफ़ी है; वह मुख्यभूमि के
हानिकर \omterm{def:b2:meiosis-heredity:vocabulary}{युग्मविकल्पी} भी उनकी मुख्यभूमि-बारंबारता पर आयात करता है, जो,
उस द्वीप के अंतःप्रजनन के नीचे, समयुग्मजों के रूप में उघड़ते हैं --- यानी वह
भार आप्रवास तथा उघड़ने के संतुलन से बैठता है।
\textbf{13.} $\hat q = \sqrt{10^{-5}} = 0.0032$; प्रभावित जन्म $\hat
q^{2} = 10^{-5}$।
\textbf{14.} $\hat q = \sqrt{10^{-5}/0.1} = 0.01$; प्रभावित जन्म
$10^{-4}$: यानी दस गुना कमज़ोर वरण उस \omterm{def:b2:meiosis-heredity:vocabulary}{युग्मविकल्पी} को तीन
गुना आम और उस रोग को दस गुना होने देता है।
\textbf{15.} $\hat p = \mu/(hs) = 10^{-5}/0.05 = 2\times 10^{-4}$;
वाहक $2\hat p = 4\times 10^{-4}$।
\textbf{16.} $\hat q = 0.0032$ से $0.01$ तक: वह \omterm{def:b2:meiosis-heredity:vocabulary}{युग्मविकल्पी} ऐसी किसी
दर से चढ़ता है जो \omterm{def:b2:genome-diversification:mutation}{उत्परिवर्तन} बाँधता है, यानी प्रति पीढ़ी $\mu = 10^{-5}$, इसलिए वह
पहुँच $\Delta q/\mu \approx 700$ पीढ़ियों की कोटि में लेती है ---
यानी \num{20000} वर्ष: और अभी लिया गया कोई भी निर्णय ऐसी समष्टि
भोगेगी जो हम जैसी बिलकुल नहीं होगी।
\textbf{17.} हर स्थल प्रति व्यक्ति $2\hat q = 0.0063$ घातक प्रतियाँ
जोड़ता है; और \num{20000} स्थलों पर, लगभग $130$ --- पर यह मानकर कि
हर जीन किसी \omterm{def:b2:meiosis-heredity:vocabulary}{अप्रभावी} घातक में उत्परिवर्तित हो सकता है; और कुछ हज़ार
अनिवार्य जीनों के चिरपरिचित आकलन के साथ, कोई व्यक्ति \omterm{def:b2:meiosis-heredity:vocabulary}{विषमयुग्मजी} दशा में
कुछ से कुछ दर्जन घातक तुल्यांक ढोता है।
\textbf{18.} हर एक किसी भिन्न स्थल पर है और हर एक दुर्लभ है, इसलिए यह
संभावना छोटी है कि कोई यादृच्छिक युगल दोनों वही एक ढोएँ;
चचेरे अपने जीनों का आठवाँ भाग वंश से साझा करते हैं, इसलिए किसी
चचेरे के कई घातकों में से हर एक के लिए यह जोखिम कि वह बच्चा \omterm{def:b2:meiosis-heredity:vocabulary}{समयुग्मजी} हो
प्रति साझा घातक $1/16$ है --- और चचेरों के विवाहों की अतिरिक्त
शिशु-मृत्यु कुछ प्रतिशत मापी गई है।
\textbf{19.} $S = 1.16\times \qty{800}{L} = \qty{930}{L}$; $R = 0.4\times
930 = \qty{370}{L}$।
\textbf{20.} पाँच पीढ़ियाँ: यानी लगभग \qty{1860}{L}; व्यवहार में छोटा,
क्योंकि वह योगात्मक प्रसरण चुक जाता है, क्योंकि $h^{2}$
मूल झुंड में आँकी गई थी, और क्योंकि वह पर्यावरण
(चारा, आवास) जिसने उन शीर्ष गायों को बनाया संचरित नहीं होता।
\textbf{21.} औसत अंतर $(1.16 + 2.4)/2 = 1.78$ मानक
विचलन, $S = \qty{1424}{L}$, $R = 0.4\times 1424 = \qty{570}{L}$।
\textbf{22.} \omterm{def:b2:asexual-reproduction:clone}{क्लोनों} में कोई आनुवंशिक प्रसरण नहीं है, इसलिए $h^{2} = 0$ और $R =
0$, $S$ जो भी हो: वंशागतिता किसी समष्टि का गुण है, यानी
उसके प्रसरण का वह हिस्सा जो आनुवंशिक है, न कि उस लक्षण का गुण।
\textbf{23.} $R = 0.7\times 0.6 = \qty{0.42}{mm}$ गहरी।
\textbf{24.} वही अनुक्रिया $\qty{370}{L}$ पाने को $S = 370/0.1 =
\qty{3700}{L}$ चाहिए, यानी $4.6$ मानक विचलन --- यानी दस हज़ार में
एक से भी कम गाय रखना: अव्यावहारिक, और इसीलिए
नीची वंशागतिता वाले लक्षण धीरे या दूसरे साधनों से सुधारे जाते हैं
(संतति-परीक्षण, बड़े परिवार)।
\textbf{25.} पतंगा: $s \approx 0.2$ से $0.3$, $0.3$ पर 28 पीढ़ियाँ
और $0.2$ पर 45, देखी गई 50 के मुक़ाबले; द्वीप: $H_{100} \approx 0.11$,
यानी मुख्यभूमि का \qty{27}{\%}; झुंड: केवल गायों से प्रति पीढ़ी \qty{370}{L},
और वरे गए साँड़ों के साथ \qty{570}{L}।
\end{solution}
